\documentclass[10pt]{article}

% --------- FORMAT PAGE ----------
\usepackage[a4paper,margin=1.5cm]{geometry}
\setlength{\columnsep}{0.5cm}
\usepackage{multicol}
\usepackage{stmaryrd}
% --------- LANGUE ----------
\usepackage[french]{babel}
\usepackage[utf8]{inputenc}
\usepackage[T1]{fontenc}


% --------- MATHS ----------
\usepackage{amsmath, amssymb, amsfonts}

% --------- COULEURS + BOX ----------
\usepackage[most]{tcolorbox}
\usepackage{xcolor}
\usepackage{newpxtext} % Pour le texte (Palatino)
\usepackage{eulervm}   % Pour les mathématiques (Euler)
\usepackage{blindtext}
\usepackage{titlesec}
\title{Exercices classiques, Maths SUP : Énoncé}
\author{Lucas 2}
\date{ }

\begin{document}
\renewcommand*\contentsname{Sommaire}
\maketitle
\setcounter{tocdepth}{3}
\tableofcontents

\section{Matrices}
\subsection{Exercices méthodes :}
\subsubsection{Diagonalisation 1.}
\subsubsection{Diagonalisation 2.}
\subsubsection{Diagonalisation 3.}

\subsubsection{Utiliser la Trace 1.}
\subsubsection{Utiliser la Trace 2.}
\subsubsection{Utiliser la Trace 3.}


\subsection{Exercices classiques.}
Soit $A,B \in \mathcal{M}_n(\mathbb{R})$, et $\lambda_0,\dots,\lambda_n \in \mathbb{K}$
telles 
\subsubsection{Nilpotence implique trace nulle.}
\subsubsection{Lemme de Hadamard.}
\subsubsection{Conditions pour que $rg(f)\leq rg(g)$, (\textit{X})}
\textbf{\underline{Énoncé :}}
Soit $E,F$ deux $\mathbb{K}$-espaces vectoriels non nuls de dimension finie, 
$f$ et $g$ dans $\mathcal{L}(E,F)$. 

Montrer que $rg(g)\leq rg(f)$ si et seulement s'il existe 
$h \in GL(F)$ et $k \in \mathcal{L}(E)$ tels que $h \circ g = f \circ k$

\subsubsection{Fonctions multiplicatives et inversibilité, (\textit{X})}
\textbf{\underline{Énoncé :}} 
Soit $f \colon \mathcal{M}_n(\mathbb{K})\rightarrow \mathbb{K}$ non constante 
telle que : $\forall A,B \in \mathcal{M}_n(\mathbb{K}), f(AB)=f(A)f(B)$.

Montrer que $f(A)=0$ si et seulement si $A$ n'est pas inversible.

\section{Polynomes}
\subsection{$P(\mathbb{U})\subset \mathbb{U}$}
\subsection{$P(\mathbb{R}) \subset \mathbb{C}$}
\subsection{Polynômes de $\mathbb{C}[X]$ induisant une surjection de $\mathbb{Q}$ dans $\mathbb{Q}$ $P(\mathbb{Q}) = \mathbb{Q}$}

\section{Espaces vectoriels, dimensions et applications linéaires}
\subsection{Première caractérisation des homothéties.}
\subsection{Seconde caractérisation des homothéties.}
\subsection{Deux preuves de la \underline{majoration de l'indice de nilpotence} en dimension finie.}
\subsection{Union de sous-espaces vectoriels stricts et supplémentaire communs.}
\subsection{Une preuve de la formule d'inversion de Pascal-Möbius.}

\section{Polynomes}
This is the first section.
\addcontentsline{toc}{section}{Unnumbered Section}
\section*{Unnumbered Section}
\subsection{Test}


\section{Second Section}

\blindtext
\end{document}